\section{Proposed Scheme}\label{sec:method}

% PA = Aim
\subsection{Design Aim}
Given a natural-language query that names a vegetable $v$, a store or market
$s$, and the price $P_t$ the user observes on day $t$, FMPA returns, for a
horizon $h\in\{7,14\}$ calendar days, a forecast price $\hat{P}_{t+h}$, the
predicted direction of change (down, flat, or up), and a short explanation of
the evidence. When the evidence is insufficient, FMPA returns a clarification
request or an explicit error instead of a number.

% PB = Based on (domain mindsets), PC = Cure (why the design helps)
\subsection{Design Principles from Domain Knowledge}
The design is based on five domain mindsets (M1--M5) that we distilled from
the literature in Section~\ref{sec:related} and from our data; each maps to a
system component.
\begin{itemize}
\item \textbf{M1 Price comparability first.} Prices are only comparable after
units and specifications (e.g., per bundle vs.\ per catty) are normalized
$\rightarrow$ a \texttt{normalize\_spec} tool.
\item \textbf{M2 Temporal integrity.} A target at $t+h$ may be used for
training only after it has matured; in our data, 1{,}386 immature samples were
conservatively removed $\rightarrow$ a label-maturity filter in the training
pipeline.
\item \textbf{M3 Simple baselines first.} Because most retail prices do not
change within a week~\cite{nakamura2008five}, every forecast is reported next
to no-change and momentum baselines $\rightarrow$ baselines are computed by the
ForecastEngine.
\item \textbf{M4 Public $\neq$ private.} Public market averages describe a
trend, not the price in a given store~\cite{meyer2004asymmetric}
$\rightarrow$ store-price anchoring \eqref{eq:modeA}.
\item \textbf{M5 Error $\neq$ utility.} A small average error does not imply a
better purchase decision $\rightarrow$ the output shows direction and evidence,
not only a point value.
\end{itemize}

% PD = Design, PE = Paper elements, PG = Graph
\subsection{System Architecture}
Fig.~\ref{fig:architecture} shows the five modules of FMPA.

\cnum{1}~\textit{User entry.} A web prototype and chat interface receive the
query. An LLM extracts the intent $(v, s, P_t, h, \text{mode})$, selects tools,
and finally explains the returned results in plain language. It is the only
module that uses the LLM.

\cnum{2}~\textit{MCP tools.} Typed tools exposed through the Model Context
Protocol~\cite{anthropic2024mcp}---\texttt{get\_history},
\texttt{normalize\_spec}, and \texttt{forecast}---validate their arguments and
return JSON. The LLM can obtain a number only through these tools.

\cnum{3}~\textit{Oracle host API.} A REST service on a cloud host serves the
tools, keeping data access and model inference outside the LLM context.

\cnum{4}~\textit{ForecastEngine.} A RandomForest~\cite{breiman2001random}
model for $h=7$ and a HistGradientBoosting model~\cite{ke2017lightgbm} for
$h=14$, both implemented with scikit-learn~\cite{pedregosa2011scikit},
estimate the relative change; the engine then applies mode A or B and computes
the baselines.

\cnum{5}~\textit{Public historical data.} Daily price records from 13 public
markets in Taichung City are used for features, training, and evaluation.

\begin{figure}[t]
\centering
% Fig. 1 -- FMPA system architecture (modules 1-5). Single-column TikZ figure.
% Module numbers must match the description in sections/03_proposed_scheme.tex.
\begin{tikzpicture}[
  font=\scriptsize,
  mod/.style={draw, rounded corners=2pt, align=center, text width=#1,
              minimum height=9mm, fill=blue!6},
  mod/.default=6.4cm,
  llm/.style={mod=#1, fill=green!10},
  llm/.default=6.4cm,
  num/.style={circle, draw, fill=white, inner sep=0.6pt, minimum size=1.15em,
              font=\scriptsize\bfseries},
  biarr/.style={{Stealth[length=1.8mm]}-{Stealth[length=1.8mm]}, semithick}
]
\node[llm] (entry) {\textbf{User entry} (web prototype / chat)\\
  LLM: intent parsing, tool selection, explanation of results};
\node[mod, below=5mm of entry] (mcp) {\textbf{MCP tools}\\
  \texttt{get\_history}\quad\texttt{normalize\_spec}\quad\texttt{forecast}};
\node[mod, below=5mm of mcp] (api) {\textbf{Oracle host API} (REST services)};
\node[mod=2.95cm, minimum height=13mm, anchor=north west] (engine)
  at ([yshift=-5mm]api.south west) {\textbf{ForecastEngine}\\
  RF (7 d) / HGB (14 d)\\ modes A/B, baselines};
\node[mod=2.95cm, minimum height=13mm, anchor=north east] (data)
  at ([yshift=-5mm]api.south east) {\textbf{Public historical data}\\
  13 Taichung public markets};

\node[num] at (entry.north west) {1};
\node[num] at (mcp.north west) {2};
\node[num] at (api.north west) {3};
\node[num] at (engine.north west) {4};
\node[num] at (data.north west) {5};

\draw[biarr] (entry) -- node[right, font=\tiny]{tool calls / JSON results} (mcp);
\draw[biarr] (mcp) -- (api);
\draw[biarr] (api.south -| engine.north) -- (engine.north);
\draw[biarr] (api.south -| data.north) -- (data.north);

\node[draw, dashed, rounded corners=3pt, inner sep=2.5mm,
      fit=(mcp)(api)(engine)(data),
      label={[font=\tiny\itshape, anchor=north, align=center, text width=6.4cm]%
      south:{deterministic tools --- every price is retrieved or computed here,
      never by the LLM}}] {};
\end{tikzpicture}

\caption{FMPA architecture. The LLM (green) only parses intent and explains
results; modules \cnum{2}--\cnum{5} inside the dashed box retrieve or compute
every number.}
\label{fig:architecture}
\end{figure}

% PF = Formulation
\subsection{Problem Formulation}
Table~\ref{tab:notation} lists the notation. Let $M_t$ be the public market
price of $v$ on day $t$. The learning target is the relative change of the
market price over $h$ calendar days,
\begin{equation}
\Delta_{t,h}=\frac{M_{t+h}-M_t}{M_t},
\label{eq:delta}
\end{equation}
and the ForecastEngine outputs an estimate $\hat{\Delta}_{t,h}$.

\textit{Mode A (store-price anchoring)} transfers the estimated relative change
to the price the user actually sees:
\begin{equation}
\hat{P}^{(A)}_{t+h}=P_t\,\bigl(1+\hat{\Delta}_{t,h}\bigr).
\label{eq:modeA}
\end{equation}

\textit{Mode B (weekly momentum adjustment)} adds a fixed fraction of the most
recent weekly momentum $m_t=(M_t-M_{t-7})/M_{t-7}$:
% VERIFY(author): confirm this exact form of mode B against the DIC3
% ForecastEngine code before submission.
\begin{equation}
\hat{P}^{(B)}_{t+h}=P_t\,\bigl(1+\hat{\Delta}_{t,h}+\alpha\, m_t\bigr),
\qquad \alpha=0.3 .
\label{eq:modeB}
\end{equation}
The coefficient $\alpha$ is fixed in advance and is not tuned on test data.

\begin{table}[t]
\caption{Notation}
\label{tab:notation}
\centering
\footnotesize
\begin{tabular}{@{}ll@{}}
\toprule
Symbol & Meaning \\
\midrule
$t$, $h$ & query day; horizon $h\in\{7,14\}$ calendar days \\
$P_t$ & price observed by the user in store $s$ on day $t$ \\
$M_t$ & public market price of vegetable $v$ on day $t$ \\
$\Delta_{t,h}$, $\hat{\Delta}_{t,h}$ & true and estimated relative change, Eq.~\eqref{eq:delta} \\
$m_t$ & weekly momentum $(M_t-M_{t-7})/M_{t-7}$ \\
$\alpha$ & fixed momentum weight in mode B ($0.3$) \\
$\theta$ & intent-confidence threshold of the agent \\
$\hat{P}^{(A)}_{t+h}$, $\hat{P}^{(B)}_{t+h}$ & forecasts of modes A and B \\
\bottomrule
\end{tabular}
\end{table}


% PJ = Jump to an example
\textit{Example.} A user reports green onions at NT\$380 per bundle, the
price reported in Taipei during Typhoon Podul~\cite{pts2025onion}. If the
engine returns $\hat{\Delta}_{t,7}=+2.4\%$, mode A gives
$\hat{P}^{(A)}_{t+7}=380\times1.024\approx$ NT\$389. The numbers in this
example are illustrative and are not experimental results.

% PH = How, PI = Implementation (flowchart + pseudocode)
\subsection{Processing Flow and Algorithm}
Fig.~\ref{fig:flowchart} shows the processing flow from input to output. The
check after intent parsing is the key safety step: if no matching history is
found, or the intent confidence is below $\theta$, FMPA asks the user for
clarification or reports ``no data'' rather than letting the LLM guess.
Algorithm~\ref{alg:fmpa} gives the corresponding pseudocode. Lines 2 and 13
are the only LLM calls; the forecast value is produced by tools in lines
11--12.

\begin{figure}[t]
\centering
% Fig. 2 -- FMPA processing flow: Input -> Check -> AI/Analytics -> Decision -> Output.
% The "no" branch of the check asks for missing fields or returns an explicit error.
\begin{tikzpicture}[
  font=\scriptsize,
  node distance=3.2mm,
  proc/.style={draw, rectangle, align=center, text width=3.6cm,
               minimum height=6mm, fill=blue!5},
  llm/.style={proc, fill=green!10},
  term/.style={draw, rounded corners=2.5mm, minimum width=1.6cm,
               minimum height=5mm, fill=gray!12},
  dec/.style={draw, diamond, aspect=2.6, align=center, inner sep=0.5pt,
              fill=yellow!20},
  err/.style={draw, rectangle, align=center, text width=1.9cm,
              minimum height=6mm, fill=red!8},
  arr/.style={-{Stealth[length=1.8mm]}, semithick}
]
\node[term] (start) {Start};
\node[proc, below=of start] (in) {Input: vegetable $v$, store price $P_t$,\\
  unit, date $t$; optional last-week price};
\node[llm, below=of in] (parse) {Web form or LLM (via MCP)\\ parses the query};
\node[dec, below=of parse] (chk) {Input valid and\\ data sufficient?};
\node[proc, below=of chk] (ai) {AI/Analytics: ForecastEngine\\
  RF ($h{=}7$) or HGB ($h{=}14$), data to $t-1$};
\node[proc, below=of ai] (decide) {Decision: mode A; add mode B only\\
  if a comparable same-store price exists};
\node[llm, below=of decide] (out) {Output: prices, $\pm5\%$ direction, dates,\\
  warnings; LLM explains, adds no numbers};
\node[term, below=of out] (end) {End};
\node[err, right=6mm of chk] (err) {Ask for missing fields, reject, or report
  error};

\draw[arr] (start) -- (in);
\draw[arr] (in) -- (parse);
\draw[arr] (parse) -- (chk);
\draw[arr] (chk) -- node[right]{yes} (ai);
\draw[arr] (ai) -- (decide);
\draw[arr] (decide) -- (out);
\draw[arr] (out) -- (end);
\draw[arr] (chk) -- node[above]{no} (err);
\draw[arr] (err.north) |- node[pos=0.25, right]{retry} (in.east);
\end{tikzpicture}

\caption{Processing flow of FMPA: Input $\rightarrow$ Check $\rightarrow$
AI/Analytics $\rightarrow$ Decision $\rightarrow$ Output, with a
clarification/error branch.}
\label{fig:flowchart}
\end{figure}

\begin{algorithm}[t]
\caption{FMPA query handling}
\label{alg:fmpa}
\begin{algorithmic}[1]
\Require user query $q$ (text), threshold $\theta$, mode $\in\{A,B\}$
\Ensure response $R$ (forecast, direction, explanation) or an error
\State $u \gets \Call{Preprocess}{q}$
\State $(v, s, P_t, h) \gets \Call{LLM.ParseIntent}{u}$
\While{$\Call{Confidence}{v, s, P_t, h} < \theta$}
  \State ask the user for clarification; update $(v, s, P_t, h)$
\EndWhile
\State $D \gets \Call{Tool.GetHistory}{v, s, t}$ \Comment{MCP tool}
\If{$D = \emptyset$}
  \State \Return error ``no matching data''
\EndIf
\State $P_t \gets \Call{Tool.NormalizeSpec}{P_t, v}$
\State $\hat{\Delta}_{t,h} \gets \Call{Tool.Forecast}{D, h}$ \Comment{RF or HGB}
\State $\hat{P}_{t+h} \gets$ Eq.~\eqref{eq:modeA} or \eqref{eq:modeB}; compute baselines
\State $R \gets \Call{LLM.Explain}{\hat{P}_{t+h}, \hat{\Delta}_{t,h}, \text{baselines}}$
\State \Return $R$
\end{algorithmic}
\end{algorithm}

% PK = Key points, PL = Later how to simulate
\subsection{Key Points}
FMPA differs from a price chatbot in three ways: (i) the forecast is anchored
on the store price the user observes rather than on a market average;
(ii) horizons are fixed calendar days, so every forecast can later be checked
against the real price on day $t+h$; and (iii) the LLM never produces a price.
Section~\ref{sec:experiments} evaluates the forecasting core against strong
simple baselines.
